\section{Bài toán, mục tiêu và phạm vi}
\subsection{Bài toán nghiên cứu}
Data drift là thay đổi phân phối đầu vào $P(X)$ giữa reference và production. Cảnh báo drift không tự chứng minh mô hình suy giảm: thay đổi có thể xảy ra ở vùng mô hình vẫn dự đoán tốt. Ngược lại, sai số có thể tăng khi quan hệ $P(Y\mid X)$ thay đổi mà detector đầu vào không phát hiện. Do đó, giám sát drift và kiểm chứng chất lượng là hai nguồn bằng chứng bổ sung.

Ở thời điểm quyết định $t$, hệ thống quan sát batch production, prediction và xác suất của mô hình đang phục vụ; chỉ được đọc nhãn đã nhận hoặc được yêu cầu trong ngân sách. Cần lựa chọn \texttt{KEEP}, \texttt{REQUEST\_LABELS}, \texttt{HOLD} hoặc \texttt{RETRAIN}. Khi không đủ bằng chứng, kết quả là chưa xác định; không gán mặc định mô hình khỏe. Ước lượng không nhãn phụ thuộc giả định về shift và calibration, không bảo đảm nhận diện mọi suy giảm~\cite{atc,cbpe}.

\subsection{Câu hỏi nghiên cứu và đóng góp dự kiến}
\begin{enumerate}
  \item \textbf{RQ1:} Kiểm chứng bằng một mẫu nhãn thật giúp giảm bao nhiêu quyết định retraining sai hoặc bỏ sót so với chỉ dùng cảnh báo drift hay ước lượng không nhãn?
  \item \textbf{RQ2:} Với cùng giới hạn ngân sách nhãn, policy kết hợp kiểm tra định kỳ và yêu cầu nhãn theo bằng chứng có duy trì chất lượng tốt hơn policy lấy nhãn theo lịch cố định, khi dùng cùng recipe retraining không?
\end{enumerate}
Đóng góp gồm: (i) cơ chế phân bổ nhãn kiểm chứng trước quyết định retraining; (ii) protocol và benchmark tái lập đánh giá tác động lên cả vòng huấn luyện; (iii) tích hợp một luồng thử nghiệm vào repo. Không đề xuất detector, estimator hay lý thuyết thống kê mới. Thành công được đánh giá bằng kết quả có đối chứng và phân tích giới hạn, kể cả khi policy đề xuất không vượt baseline.

\subsection{Ranh giới thực hiện trong 5 tháng}
\begin{itemize}
  \item \textbf{Cốt lõi:} Hai dataset nhị phân tabular, một model family có xác suất lớp; CBPE chính, ATC đối chứng; nhãn kiểm chứng lấy ngẫu nhiên; một recipe full-retraining; đánh giá offline và demo tích hợp cục bộ.
  \item \textbf{Giả định nhãn:} Dataset thực nghiệm có nhãn đầy đủ ở phía evaluator. Hệ thống chỉ nhận phần được tiết lộ theo ngân sách sau mỗi batch. Nhãn đúng, chi phí một đơn vị/mẫu; phần còn lại chưa được phép đọc. Yêu cầu nhãn được mô phỏng, không phải nghiên cứu hành vi người gán nhãn.
  \item \textbf{Ngoài phạm vi:} Taxonomy nguyên nhân, causal diagnosis, học policy bằng reinforcement learning, tối ưu sampler, incremental training, pseudo-labeling, nhãn trễ/sửa nhãn phức tạp, human study và canary production.
\end{itemize}
Trường hợp không lấy được bất kỳ nhãn production nào chỉ dùng kiểm tra giới hạn giám sát. Không cam kết supervised retraining trên dữ liệu không nhãn. Tên dataset, cấu hình model, mức ngân sách và ngưỡng được chọn qua pilot ở tháng 1 rồi đóng băng trước test.
